\documentclass[11pt]{article}
\usepackage[margin=1in]{geometry}
\usepackage{amsmath}
\usepackage{booktabs}

\title{Standardized Logistic Regression Is a Strong Baseline on Handwritten Digits}
\author{Paper2Proof canary set\thanks{A small paper written for the Paper2Proof evaluation set (eval set A).}}
\date{2026}

\begin{document}
\maketitle

\begin{abstract}
We revisit multinomial logistic regression on the scikit-learn handwritten digits dataset.
With per-feature standardization and an L2 penalty, a plain linear model reaches a test accuracy
of 0.965 averaged over three random train/test splits, which makes it a strong and cheap
reference point for small image classification benchmarks.
\end{abstract}

\section{Introduction}
Small benchmarks are often used to compare new classifiers, but the linear baseline is rarely
tuned. We report the accuracy of a standardized logistic regression on the digits dataset under
a fully specified protocol, so that the number can be reproduced exactly.

\section{Method}
The model is a multinomial logistic regression with an L2 penalty ($C = 1.0$) trained with the
L-BFGS solver for at most 2000 iterations. Each of the 64 pixel features is standardized to zero
mean and unit variance using statistics computed on the training split only.

\section{Experimental setup}
\paragraph{Data.} The digits dataset shipped with scikit-learn (\texttt{load\_digits}): 1797
grayscale images of $8 \times 8$ pixels, 10 classes. No download is needed.
\paragraph{Splits.} For each seed $s \in \{0, 1, 2\}$ we draw a stratified train/test split with
30\,\% of the samples in the test set (\texttt{train\_test\_split(test\_size=0.3,
random\_state=s, stratify=y)}).
\paragraph{Metric.} Classification accuracy on the test split. We report the mean and the sample
standard deviation over the three seeds.
\paragraph{Software.} Python 3.12, scikit-learn 1.9.1, CPU only. One run takes well under a
second.

\section{Results}
Table~\ref{tab:digits-logreg} reports the result. Standardized logistic regression reaches a
test accuracy of 0.965 $\pm$ 0.009 on digits.

\begin{table}[t]
\centering
\caption{Test accuracy on the digits dataset (mean $\pm$ standard deviation over seeds 0, 1, 2).}
\label{tab:digits-logreg}
\begin{tabular}{lc}
\toprule
Model & Accuracy \\
\midrule
Logistic regression (standardized, $C = 1.0$) & $0.965 \pm 0.009$ \\
\bottomrule
\end{tabular}
\end{table}

\section{Conclusion}
A standardized linear model is a strong baseline on digits and should be reported alongside any
new method evaluated on this dataset.

\end{document}
